\section{Counting a rival's money}
\label{sec:bargain}

A trade changes what the other player can raise and build. This section reports how agents respond to offers and what their texts say about a counterparty's finances. No reference value exists for a trade in a four-player game, so nothing here measures whether a response was good.

\subsection{All 920 responses}
The games contain 526 negotiation threads and 920 responses to offers: 113 acceptances, 394 counters and 413 rejections. Sol and Opus mostly counter (140 of 235 and 141 of 189 responses). Gemini and Grok mostly reject (180 of 242 and 149 of 254). Each accepts between 26 and 30 offers.

Table~\ref{tab:opening} gives every opening offer by form and by whose color group the trade would complete. Opening swaps that complete only the proposer's group are rejected 95 times in 116. Swaps that complete neither side's group are rejected 37 times in 72, a difference of 30.5 percentage points.

\begin{table*}[t]\centering\small
\begin{tabular}{@{}lrrrr@{}}\toprule
Trade would newly complete a group for & Swap & Cash bid & Cash ask & Other\\\midrule
\rows{gen_opening_rows}
\bottomrule
\end{tabular}
\caption{Every opening offer where acceptance was legal: $N$ [accepted/countered/rejected]. A swap moves deeds both ways, a cash bid asks for deeds and a cash ask offers them.}\label{tab:opening}
\end{table*}

The difference depends on who proposed. Grok wrote 36 of the 72 neither-completing swaps and fifteen of their seventeen acceptances, and thirteen of those fifteen offers came from calls whose assistant content named a different action. Without Grok-authored offers the rejection counts are 90 of 108 and 24 of 36, a difference of 16.7 points. Resampling the five seed blocks gives a percentile range of $-3.8$ to $31.9$ points for that figure. The range rests on five blocks and a contrast chosen after the fact. It shows instability and does not show that the association is absent. Opening cash bids differ by $-5.0$ points when pooled and by $-0.2$ to $+8.5$ points within recipients. Proposer-completing swaps are rejected more often in this cohort, and the size of the gap follows the mix of proposers. The data do not identify an effect of completion (Appendix~\ref{app:bargain}).

A separate population bears on what refusals say. Seventy-four rejections concern offers that would have completed the recipient's own group. Two model readings made for an earlier version of this study agree that 61 of the 74 texts cite a consequence for the rival. We did not repeat that reading and no person has checked it. These are different offers from the swap contrast, so the reading cannot explain it.

\subsection{Four incidents}
The system prompt permits deception in public messages, so no conduct below breaks an instruction. We selected the incidents for clarity. They show neither how often a tactic occurs nor whether it works. Three more are in Appendix~\ref{app:incidents}.

\incident{3}{Refusing one's own monopoly (G15).} Sol offers both its greens for Grok's Ventnor and \$100, which would complete a group for each. Grok holds \$333 and Sol \$730. Grok refuses, and its note counts both budgets (p4 is Sol):
\QB{g15-refusal-note}
The arithmetic is right. Sol would hold \$830 and Grok \$233, and a first even layer of green houses costs \$600. Sol raises the offer to both greens, Kentucky and \$250, and says this leaves Grok \$583 against its own \$480. Grok answers \Q{g15-final} Eighteen turns later Grok is eliminated and its deeds pass to Gemini. The refusal is argued from the rival's building budget. Whether it helped Grok is unknown.

\incident{4}{One deed, two prices (G2).} Grok owns no yellow deed and bids \$200 for Sol's Marvin Gardens to keep it from Opus. Sol asks \$450 and Grok objects \Q{g2-blocker-price} Opus, who owns the other two yellows and \$845, then bids \$325. Sol asks \$650:
\QB{g2-price}
Sol's note explains the price with a figure that is wrong, where p4 is Opus: \Q{g2-note} A yellow house costs \$150, so \$520 buys three. The pricing rule is clear and the rival's budget is overstated threefold. Opus calls \$450 final in public while its note sets a ceiling of about \$475. Sol keeps the deed.

\incident{5}{A bid timed to the seller's cash (G7).} Sol buys Park Place, builds three houses and has \$84 left. Gemini's note sees an opening (p4 is Sol):
\QB{g7-timed}
Gemini bids \$350 for New York Avenue, raises to \$520, and tells Sol the sum would give it exactly the \$600 that three more houses cost. Sol holds at \$700 and tells Gemini it can mortgage its other deeds to fund building. Gemini pays \$700, mortgages five deeds and builds four orange houses in the same turn. Each side computed the other's building budget. Earlier in that game Grok told Sol \Q{g7-public} while its note read \Q{g7-private}. Sol took \$600.

\incident{6}{Financing a rival to keep an estate from a creditor (G13).} Grok has \$35 and could owe Sol a \$200 rent on its next roll. It offers Opus a mortgaged deed for \$200, calling Sol ``gpt'':
\QB{g13-pitch}
Opus counters at \$130. Grok's note works out that \Q{g13-floor} Its provider field adds \Q{g13-omit} Opus holds \$1,926 and pays the \$200. Its provider field says that \Q{g13-buyer} Player 2 is Sol. The buyer paid to keep a third player from inheriting a color group. The seller's arithmetic was correct, and it left one financing option out of the pitch. Grok was eliminated on a \$200 debt to Sol 43 turns later.
